\section{Exact searches on incidence graphs}
The searches in this section test concrete graphs below their star-circuit sizes. Every reported search stopped without a satisfying circuit or an unsatisfiability proof, so none supplies an additional lower bound.

For $\operatorname{PG}(2,3)$ there are thirteen points and thirteen lines, each of degree four. Every cover therefore has at least $52/4=13$ vertices, and either bipartition class attains thirteen. The star count is sixty-four. For the full affine plane of order three, the nine points have degree four and the twelve lines degree three; the same counting argument gives cover number nine and star count forty-four.

The encoder chooses a topological list of binary gates, their AND/OR types and their two earlier input wires. Correctness is imposed on every edge assignment and every maximal independent set. This is sufficient: every positive input contains an edge, and every negative input extends to a maximal independent set, so monotonicity determines all other values. Input permutations by graph automorphisms reduce the first gate's two inputs to orbit representatives; the rest of the circuit is unrestricted by symmetry.

The one-sided theorem removes circuits with marked degree one, and the three-missing-witness theorem removes circuits having fewer than four missing quadratic monomials. Formal degree-one and degree-two occurrence bits track these conditions without expanding the full polynomial. The terminal-product and common- or private-mask exclusions of \citet[Sections~9--10]{OneSidedCompanion} are applied to recognized circuit configurations as additional pruning. There was no candidate requiring these later lazy exclusions in the reported attempts.

\begin{center}
\small
\begin{tabular}{@{}lrrrr@{}}
\toprule
Graph & Gates tested & Negative tests & Symmetries used & Result\\\midrule
Full affine plane of order $3$ &43&197&432&undecided\\
Projective plane of order $3$ &63&704&78&undecided\\
Subdivided Petersen graph &38&579&120&undecided\\
Tutte--Coxeter graph &58&3264&720&undecided\\
\bottomrule
\end{tabular}
\end{center}
Each attempt allowed three hundred seconds of solver time after encoding. The corresponding CNF sizes, in variables and clauses, were
\[
\begin{array}{c|r@{\ ,\ }r}
\text{affine}&54\,910&1\,722\,686\\
\text{projective}&193\,413&9\,748\,983\\
\text{subdivided Petersen}&90\,164&3\,347\,138\\
\text{Tutte--Coxeter}&619\,894&35\,501\,275.
\end{array}
\]
The two latter star counts are thirty-nine and fifty-nine. The original Petersen graph has two disjoint five-cycles covering its ten vertices. Assign to each vertex the next edge on its cycle; the ten distinct assigned edges become distinct neighbors after subdivision. This gives a matching saturating the original vertices, so the subdivided graph has cover number ten and its thirty edges give the star count thirty-nine. The cubic bipartite Tutte--Coxeter graph has fifteen vertices on each side; the regular bipartite incidence count gives cover number fifteen and star count fifty-nine.

As a calibration, fixing the explicit forty-four-gate affine star circuit made the encoding satisfiable, and its decoded gate list was independently evaluated on every critical assignment. This checks a positive example; it does not certify the search encodings as infeasibility proofs. The useful-gate convention tests exactly the displayed gate count, so even an eventual unsatisfiability result would need a separate justification before it covered every smaller count. The program \emph{Incidence-circuit synthesis}, its commands, encodings and results accompany the workbench.
